\documentclass[10pt]{beamer}
\usetheme{metropolis}
\usepackage{graphicx}
\usepackage{listings}
\usepackage{booktabs}
\usepackage{xcolor}
\definecolor{codebg}{HTML}{F5F5F5}
\definecolor{codegreen}{HTML}{2E7D32}
\definecolor{codegray}{HTML}{757575}
\definecolor{codeblue}{HTML}{1565C0}
\lstdefinestyle{rstyle}{language=R,backgroundcolor=\color{codebg},basicstyle=\ttfamily\scriptsize,
  keywordstyle=\color{codeblue}\bfseries,stringstyle=\color{codegreen},
  commentstyle=\color{codegray}\itshape,breaklines=true,frame=single,
  rulecolor=\color{codegray!30},showstringspaces=false,tabsize=2,
  morekeywords={library,ggplot,aes,geom_point,geom_smooth,geom_hex,geom_density_2d,
    geom_density_2d_filled,geom_rug,scale_x_log10,scale_color_manual,scale_size,
    ggMarginal,theme_minimal,labs,coord_cartesian}}
\lstdefinestyle{pythonstyle}{language=Python,backgroundcolor=\color{codebg},basicstyle=\ttfamily\scriptsize,
  keywordstyle=\color{codeblue}\bfseries,stringstyle=\color{codegreen},
  commentstyle=\color{codegray}\itshape,breaklines=true,frame=single,
  rulecolor=\color{codegray!30},showstringspaces=false,tabsize=4,
  morekeywords={import,from,as,pd,plt,sns,np,fig,ax,scatter,hexbin,contourf,
    set_xscale,polyfit,plot,fill_between,GridSpec,jointplot,lmplot,regplot}}
\title{Relationships: Scatter, Bubble \& Hexbin}
\subtitle{Workshop 3 --- Module 4}
\author{Prof.Asoc.\ Endri Raco}
\institute{Data Visualization for Data Scientists \\ UNYT Tirana}
\date{2025/26}
\begin{document}

\begin{frame}
  \titlepage
\end{frame}

\begin{frame}{Module Roadmap}
  \begin{enumerate}
    \item Scatter plot anatomy: position, colour, size, regression
    \item Four correlation patterns and Anscombe's warning
    \item Overplotting: alpha, jitter, hexbin, 2D density
    \item Bubble charts: encoding a 3rd and 4th variable
    \item Smoothers: linear, LOESS, and when to use which
    \item Marginal distributions: scatter + histogram/rug
  \end{enumerate}
  \begin{exampleblock}{Learning Outcome}
    You will diagnose and fix overplotting, choose the right smoother,
    and produce Gapminder-style bubble charts encoding up to five variables.
  \end{exampleblock}
\end{frame}

\section{Scatter Plot Anatomy}

\begin{frame}{Position + Colour + Regression + Confidence Band}
  \begin{center}
    \includegraphics[width=0.78\textwidth]{figures_w03m04/scatter_anatomy}
  \end{center}
\end{frame}

\begin{frame}{Four Correlation Patterns}
  \begin{center}
    \includegraphics[width=0.95\textwidth]{figures_w03m04/correlation_patterns}
  \end{center}
  \begin{alertblock}{Pro-Tip}
    Pearson's $r$ misses non-linear relationships (panel 4).
    \textbf{Always plot the scatter} before trusting a correlation
    coefficient --- this is Anscombe's lesson from W01-M01.
  \end{alertblock}
\end{frame}

\begin{frame}[fragile]{Scatter in R vs Python}
  \begin{columns}[T]
    \begin{column}{0.48\textwidth}
      \textbf{R (ggplot2)}
\begin{lstlisting}[style=rstyle]
ggplot(apps, aes(
    x = Reviews, y = Rating,
    color = Type)) +
  geom_point(alpha = 0.3,
    size = 0.8) +
  geom_smooth(method = "lm",
    se = TRUE, linewidth = 0.8) +
  scale_x_log10(
    labels = scales::comma) +
  scale_color_manual(values =
    c(Free = "#1565C0",
      Paid = "#E53935")) +
  theme_minimal() +
  labs(x = "Reviews (log)",
       y = "Rating")
\end{lstlisting}
    \end{column}
    \begin{column}{0.48\textwidth}
      \textbf{Python}
\begin{lstlisting}[style=pythonstyle]
# seaborn
sns.scatterplot(data=apps,
    x="Reviews", y="Rating",
    hue="Type", alpha=0.3, s=10,
    palette={"Free":"#1565C0",
             "Paid":"#E53935"})
ax.set_xscale("log")

# + regression
sns.lmplot(data=apps,
    x="Reviews", y="Rating",
    hue="Type", ci=95,
    height=5, aspect=1.3,
    scatter_kws={"alpha":0.3,"s":8})

# matplotlib manual
ax.scatter(x, y, s=8, c=color,
    alpha=0.3, edgecolors="none")
m, b = np.polyfit(x, y, 1)
ax.plot(xs, m*xs+b, c="grey")
\end{lstlisting}
    \end{column}
  \end{columns}
\end{frame}

\section{Overplotting}

\begin{frame}{The Overplotting Cascade}
  \begin{center}
    \includegraphics[width=0.95\textwidth]{figures_w03m04/overplotting}
  \end{center}
  \begin{exampleblock}{Data Insight}
    As $n$ grows: \textbf{alpha} works to $\sim$2K; \textbf{hexbin}
    works to $\sim$50K; \textbf{2D density contour} works for any $n$.
    The progression alpha $\rightarrow$ hexbin $\rightarrow$ density
    is a scale-adaptive overplotting solution.
  \end{exampleblock}
\end{frame}

\begin{frame}[fragile]{Overplotting Solutions in Code}
  \begin{columns}[T]
    \begin{column}{0.48\textwidth}
      \textbf{R}
\begin{lstlisting}[style=rstyle]
# Alpha transparency
geom_point(alpha = 0.05, size = 0.5)

# Hexbin
geom_hex(bins = 30) +
  scale_fill_viridis_c()

# 2D density contour
geom_density_2d(color = "grey50")

# 2D density filled
geom_density_2d_filled(alpha = 0.7)
\end{lstlisting}
    \end{column}
    \begin{column}{0.48\textwidth}
      \textbf{Python}
\begin{lstlisting}[style=pythonstyle]
# Alpha
ax.scatter(x, y, alpha=0.05, s=3)

# Hexbin
ax.hexbin(x, y, gridsize=30,
    cmap="YlGnBu", mincnt=1)

# 2D KDE contour
from scipy.stats import gaussian_kde
kde = gaussian_kde(np.vstack([x,y]))
ax.contourf(xg, yg,
    kde(pts).reshape(xg.shape),
    levels=10, cmap="YlGnBu")

# seaborn jointplot
sns.jointplot(data=df, x="x", y="y",
    kind="hex")  # or "kde"
\end{lstlisting}
    \end{column}
  \end{columns}
\end{frame}

\begin{frame}{Hexbin: Density-Encoded Scatter}
  \begin{center}
    \includegraphics[width=0.72\textwidth]{figures_w03m04/hexbin}
  \end{center}
\end{frame}

\section{Bubble Charts}

\begin{frame}{Bubble: 4 Variables in One Chart}
  \begin{center}
    \includegraphics[width=0.82\textwidth]{figures_w03m04/bubble}
  \end{center}
\end{frame}

\begin{frame}[fragile]{Bubble Chart in Code}
  \begin{columns}[T]
    \begin{column}{0.48\textwidth}
      \textbf{R}
\begin{lstlisting}[style=rstyle]
ggplot(gap, aes(
    x = gdpPercap,
    y = lifeExp,
    color = continent,
    size = pop)) +
  geom_point(alpha = 0.5) +
  scale_x_log10(
    labels = scales::dollar) +
  scale_size(
    range = c(1, 15),
    guide = "none") +
  scale_color_brewer(
    palette = "Set2") +
  theme_minimal() +
  labs(title = "Gapminder",
    x = "GDP per Capita ($, log)",
    y = "Life Expectancy")
\end{lstlisting}
    \end{column}
    \begin{column}{0.48\textwidth}
      \textbf{Python}
\begin{lstlisting}[style=pythonstyle]
for cont, c in palette.items():
    m = df["continent"] == cont
    ax.scatter(
        df.loc[m, "gdp"],
        df.loc[m, "life_exp"],
        s=df.loc[m, "pop"] / scale,
        c=c, alpha=0.45,
        edgecolors="white", lw=0.5,
        label=cont)

ax.set_xscale("log")
ax.legend(fontsize=7,
    markerscale=0.5)

# Or plotly for interactive:
# px.scatter(df, x="gdp",
#     y="lifeExp", size="pop",
#     color="continent",
#     hover_name="country")
\end{lstlisting}
    \end{column}
  \end{columns}
  \begin{alertblock}{Pro-Tip}
    Map size to \textbf{area}, not radius: \texttt{scale\_size()} in R
    and \texttt{s=pop/scale} in Python both use area by default. Cap at
    $\leq$5 variables total (x, y, color, size, shape) to avoid
    cognitive overload.
  \end{alertblock}
\end{frame}

\section{Smoothers}

\begin{frame}{Linear vs LOESS: Choosing the Right Smoother}
  \begin{center}
    \includegraphics[width=0.92\textwidth]{figures_w03m04/smoothers}
  \end{center}
  \begin{exampleblock}{Data Insight}
    Use \texttt{method="lm"} when the relationship is approximately
    linear. Use \texttt{method="loess"} (default in ggplot2) or
    \texttt{method="gam"} for non-linear patterns. LOESS is local
    regression: it fits separate models in sliding windows, producing
    a flexible curve.
  \end{exampleblock}
\end{frame}

\section{Marginal Distributions}

\begin{frame}{Scatter + Marginal Histograms}
  \begin{center}
    \includegraphics[width=0.72\textwidth]{figures_w03m04/marginal}
  \end{center}
  \begin{alertblock}{Pro-Tip}
    In R: \texttt{ggExtra::ggMarginal(p, type = "histogram")}. In
    Python: \texttt{sns.jointplot(kind="scatter")} or manual
    \texttt{GridSpec} with histograms on top/right. Marginal distributions
    reveal univariate patterns (skewness, bimodality) that the scatter
    alone hides.
  \end{alertblock}
\end{frame}

\section{Complete Examples}

\begin{frame}[fragile]{R: Reviews vs Rating (Publication Quality)}
  \begin{columns}[T]
    \begin{column}{0.52\textwidth}
\begin{lstlisting}[style=rstyle]
apps_clean |>
  filter(!is.na(Rating),
         Reviews > 0) |>
  ggplot(aes(x = Reviews,
    y = Rating, color = Type)) +
  geom_point(alpha = 0.3,
    size = 0.8) +
  geom_smooth(method = "lm",
    se = TRUE, linewidth = 0.8) +
  scale_x_log10(
    labels = scales::comma) +
  scale_color_manual(values =
    c(Free = "#1565C0",
      Paid = "#E53935")) +
  coord_cartesian(
    ylim = c(1, 5)) +
  theme_minimal() +
  labs(
    title = "More Reviews
      = Higher Rating",
    subtitle = "Google Play Store",
    x = "Reviews (log)",
    y = "Rating",
    color = "Type")
\end{lstlisting}
    \end{column}
    \begin{column}{0.45\textwidth}
      \includegraphics[width=\textwidth]{figures_w03m04/r_output}
    \end{column}
  \end{columns}
\end{frame}

\begin{frame}[fragile]{Python: Same Chart}
  \begin{columns}[T]
    \begin{column}{0.52\textwidth}
\begin{lstlisting}[style=pythonstyle]
fig, ax = plt.subplots(figsize=(6,5))

for t, c in [("Free","#2E7D32"),
             ("Paid","#E53935")]:
    m = apps["Type"] == t
    ax.scatter(
        apps.loc[m, "Reviews"],
        apps.loc[m, "Rating"],
        s=10, c=c, alpha=0.35,
        edgecolors="none", label=t)

ax.set_xscale("log")
ax.set_ylim(1, 5)
ax.legend(fontsize=7, title="Type")
ax.set_xlabel("Reviews (log)")
ax.set_ylabel("Rating")
ax.set_title(
    "More Reviews = Higher Rating",
    fontweight="bold")
ax.spines["top"].set_visible(False)
ax.spines["right"].set_visible(False)
\end{lstlisting}
    \end{column}
    \begin{column}{0.45\textwidth}
      \includegraphics[width=\textwidth]{figures_w03m04/py_output}
    \end{column}
  \end{columns}
\end{frame}

\section{Synthesis}

\begin{frame}{Relationship Chart Selection Guide}
  \begin{center}
    \includegraphics[width=0.92\textwidth]{figures_w03m04/decision_card}
  \end{center}
\end{frame}

\begin{frame}{Key Takeaways}
  \begin{enumerate}
    \item \textbf{Scatter} is the default for two continuous variables.
          Add colour for a 3rd (categorical), size for a 4th (continuous).
    \item \textbf{Always plot the scatter} before trusting $r$ ---
          non-linear patterns and outliers are invisible in a number.
    \item \textbf{Overplotting cascade}: alpha ($n < 2K$) $\rightarrow$
          hexbin ($n < 50K$) $\rightarrow$ 2D density (any $n$).
    \item \textbf{Bubble} encodes 4--5 variables; cap at 5 to avoid
          cognitive overload. Size maps to area, not radius.
    \item \textbf{method="lm"} for linear; \textbf{method="loess"} for
          non-linear. LOESS is the ggplot2 default for $n < 1000$.
    \item \textbf{Marginal distributions} add univariate context:
          \texttt{ggMarginal()} or \texttt{sns.jointplot()}.
  \end{enumerate}
\end{frame}

\begin{frame}{Essential Readings}
  \begin{itemize}
    \item Cleveland, W.~S. \& McGill, R. (1984). ``Graphical Perception.''
          \emph{JASA}, 79(387), 531--554.
    \item Wickham, H. (2016). \emph{ggplot2}, Chapter 5 (Scatter + Smoothers).
    \item Wilke, C.~O. (2019). \emph{Fundamentals of Data Visualization},
          Chapters 12--13 (Scatter, Bubble).
    \item Rosling, H. et al. (2018). \emph{Factfulness}. (Gapminder
          bubble charts in context.)
  \end{itemize}
  \vspace{6pt}
  \textbf{Next Module}: Proportions --- Pie, Donut, Treemap \& Waffle (W03-M05)
\end{frame}

\end{document}
